\section*{5. Ax 的一个定理 [1] (亦见 [4], 第 7 节)}

\subsubsection*{定理 5.0.}
设 $V$ 是 $\mathbb P^{n+1}$ 中次数为 $d>0$、定义于 $\mathbb F_q$ 上的超曲面,
$j$ 为整数. 若

\begin{equation}\tag{5.0.1}
jd<n+2,
\end{equation}

则

\begin{equation}\tag{5.0.2}
Z(t,V/\mathbb F_q)\prod_{i=0}^{n+1}(1-q^it)
\in1+q^jt\mathbb Z[[q^jt]],
\end{equation}

且下式的每个零点或极点的倒数

\begin{equation}\tag{5.0.3}
Z(t,V/\mathbb F_q)\prod_{i=0}^{n+1}(1-q^it)
\end{equation}

\noindent
均可写成 $q^j$ 乘以一个代数整数.

\subsubsection*{定理 5.0 bis.}
设 $V_{\mathrm{aff}}$ 是 $\mathbb A^{n+2}$ 中由 $\mathbb F_q$ 上次数为 $d$ 的齐次式
定义的超曲面.\EditorialFootnote{法文原书第428页写作“非齐次式” (forme inhomogène), 英文底本改作齐次式. 下面 (5.0.7) 的仿射锥点数关系及两定理结论等价的证明使用齐次假设; 此处按英文保留这一证明范围, 并记录其与印本的差异.} 若

\begin{equation}\tag{5.0.4}
jd<n+2,
\end{equation}

则

\begin{equation}\tag{5.0.5}
Z(t,V_{\mathrm{aff}}/\mathbb F_q)
\in1+q^jt\mathbb Z[[q^jt]].
\end{equation}

特别地,

\begin{equation}\tag{5.0.6}
\operatorname{card}\bigl(V_{\mathrm{aff}}(\mathbb F_q)\bigr)
\equiv0\pmod{q^j},
\end{equation}

当 $j\geq1$ 时, 因而\EditorialFootnote{原书在此未重申 $j\geq1$. (5.0.6 bis) 末尾的严格正性需此条件; 定理的前述整除与 zeta 断言仍按原文允许整数 $j$.}

\begin{equation}\tag{5.0.6 bis}
\operatorname{card}\bigl(V(\mathbb F_q)\bigr)
\geq\frac{q^j-1}{q-1}>0,
\end{equation}

\noindent
这加强了 Chevalley--Warning 定理~[8].

\medskip
\noindent\emph{底本注.}
印本将定理~5.0 bis 的假设编号为 (5.0.7), 与随后点数恒等式的编号重复.
校订法文层和英文层采用连续编号 (5.0.4); 照录法文层保留印本编号.

\medskip
\noindent\emph{证明.}
在 zeta 函数中取 $t$ 的系数, 即知 (5.0.5) $\Rightarrow$ (5.0.6).
断言 (5.0.5) 与 (5.0.2) 等价,\EditorialFootnote{原书第429页及英文底本写作 (5.0.4), 但 (5.0.4) 是两定理共有的次数假设; 这里比较的是结论 (5.0.5) 与 (5.0.2).} 因为

\begin{equation}\tag{5.0.7}
\operatorname{card}V_{\mathrm{aff}}(\mathbb F_{q^\nu})
=1+(q^\nu-1)\operatorname{card}V(\mathbb F_{q^\nu}),
\end{equation}

从而

\begin{equation}\tag{5.0.8}
Z(t,V_{\mathrm{aff}}/\mathbb F_q)
=\frac{Z(qt,V/\mathbb F_q)}
{Z(t,V/\mathbb F_q)(1-t)}
\end{equation}

且

\begin{equation}\tag{5.0.9}
Z(t,V/\mathbb F_q)
=\prod_{i\geq0}
\left[(1-q^it)Z(q^it,V_{\mathrm{aff}}/\mathbb F_q)\right]^{-1}.
\end{equation}

\begin{editorialnote}
原书第429页及英文底本在 (5.0.8) 中颠倒了 $Z(t,V)$ 与 $Z(qt,V)$, 并在 (5.0.9) 中漏去各因子的逆.
由 (5.0.7) 及定义 (4.5.31),
$\log Z(t,V_{\mathrm{aff}})=-\log(1-t)+\log Z(qt,V)-\log Z(t,V)$,
故得到上列两式. 例如 $V$ 为一个有理点时, 它们给出 $Z(t,V_{\mathrm{aff}})=(1-qt)^{-1}$.
\end{editorialnote}

\noindent
由 Fatou 定理~[3] 以及 zeta 函数的有理性, 断言 (5.0.2) 和 (5.0.3) 同样等价.
由于

\begin{equation}\tag{5.0.10}
Z(t,V/\mathbb F_q)\in\mathbb Z[[t]],
\end{equation}

只需证明

\begin{equation}\tag{5.0.11}
Z(t,V/\mathbb F_q)\prod_{i=0}^{n+1}(1-q^it)
\in F^j(1+t\mathcal O[[t]]).
\end{equation}

根据 (4.5.33) 和 (4.6.4)--(4.6.5), 只需对
$\varnothing\ne B\subset\{1,\ldots,n+2\}$ 证明

\begin{equation}\tag{5.0.12}
\Delta_B(t)\in F^{j+\operatorname{card}(B)-(n+2)}.
\end{equation}

只要证明下式, 即可应用引理~4.6.6 完成证明

\begin{equation}\tag{5.0.13}
jd<n+2
\quad\Longrightarrow\quad
\bigl(j+\operatorname{card}(B)-(n+2)\bigr)d
<\operatorname{card}(B).
\end{equation}

这可改写为

\begin{equation}\tag{5.0.14}
jd<n+2
\quad\Longrightarrow\quad
(j-(n+2))d<\operatorname{card}(B)(1-d).
\end{equation}

但

\[
\inf_{\varnothing\ne B\subset\{1,\ldots,n+2\}}
\operatorname{card}(B)(1-d)
=(n+2)(1-d),
\]

故只需证明

\begin{equation}\tag{5.0.15}
jd<n+2
\quad\Longrightarrow\quad
(j-(n+2))d<(n+2)(1-d),
\end{equation}

此式成立.

\medskip
\noindent\emph{底本注.}
在 (5.0.9) 中, 印本将因子 $(1-q^it)$ 误印为不可能的 $(1-q^if)$,
且末段引用引理~4.4.6, 而非实际使用的整性引理 4.6.6.
上文均已校正; 照录法文层保留印本形式.

